\documentclass[14pt,a4paper,oneside]{extbook}
\usepackage[utf8]{vietnam}
\usepackage{amsmath,amssymb,mathrsfs,fancyhdr,enumerate,multirow,makecell}
\usepackage{tkz-euclide,tikz-3dplot,tikz,tkz-tab}
\usetikzlibrary{arrows.meta,calc,intersections,angles,quotes,backgrounds,shapes.geometric,patterns,through,shadows,positioning}
\usepackage{yhmath}
\usepackage[top=0.7cm, bottom=0.7cm, left=1.0cm, right=1.0cm, headheight=16pt, headsep=3mm, footskip=5mm]{geometry}
\usepackage[hidelinks,unicode]{hyperref}
\usepackage{qrcode}
\usepackage[many]{tcolorbox}
\usepackage{fontawesome}

\pagestyle{fancy}
\fancyhf{}
\fancyhead[L]{\footnotesize\bfseries\color{teal!85!black}\faBook~SỔ TAY KIẾN THỨC TOÁN 4}
\fancyhead[R]{\footnotesize\bfseries\color{blue!85!black}Lớp Toán Thầy Lĩnh -- loptoanthaylinh.vn}
\fancyfoot[C]{\footnotesize\bfseries\color{gray!80!black}-- Sổ tay kiến thức: Bài 57. Quy đồng mẫu số các phân số --}
\renewcommand{\headrulewidth}{0.4pt}

\begin{document}

% --- BANNER TIÊU ĐỀ BÀI HỌC ---
\noindent
\begin{tcolorbox}[
  enhanced,
  colback=teal!8!white,
  colframe=teal!80!black,
  arc=3mm, boxrule=1pt,
  left=3mm, right=3mm, top=2mm, bottom=2mm,
  halign=center
]
  {\Large\bfseries\color{teal!90!black}\faBookmark\ BÀI 57: QUY ĐỒNG MẪU SỐ CÁC PHÂN SỐ}\\[1.5pt]
  {\footnotesize\itshape\color{black!75} Mẫu số chung, phương pháp nhân chéo và trường hợp mẫu lớn chia hết cho mẫu bé trên 1 trang A4}
\end{tcolorbox}

\vspace{-2.5mm}

\noindent
\begin{minipage}[t]{0.49\textwidth}\vspace{0pt}
\begin{tcolorbox}[
  enhanced,
  colback=blue!4!white,
  colframe=blue!75!black,
  arc=2mm, boxrule=0.8pt,
  left=2mm, right=2mm, top=1.8mm, bottom=1.8mm,
  title={\footnotesize\bfseries\color{white}\faArrowsH\ 1. Trường hợp chung (Nhân chéo)}
]
{\footnotesize
Quy đồng mẫu số $\dfrac{a}{b}$ và $\dfrac{c}{d}$:\\
\textbf{Mẫu số chung (MSC) = $b \times d$}:
\begin{itemize}\setlength{\itemsep}{1pt}\setlength{\leftskip}{-2.5mm}
  \item Lấy cả tử và mẫu của $\dfrac{a}{b}$ nhân với mẫu số $d$:
  $$\dfrac{a}{b} = \dfrac{a \times d}{b \times d}$$
  \item Lấy cả tử và mẫu của $\dfrac{c}{d}$ nhân với mẫu số $b$:
  $$\dfrac{c}{d} = \dfrac{c \times b}{d \times b}$$
\end{itemize}
}
\end{tcolorbox}
\end{minipage}
\hfill
\begin{minipage}[t]{0.49\textwidth}\vspace{0pt}
\begin{tcolorbox}[
  enhanced,
  colback=teal!4!white,
  colframe=teal!75!black,
  arc=2mm, boxrule=0.8pt,
  left=2mm, right=2mm, top=1.8mm, bottom=1.8mm,
  title={\footnotesize\bfseries\color{white}\faLightbulbO\ 2. Mẫu lớn chia hết cho mẫu bé}
]
{\footnotesize
Khi một mẫu số chia hết cho mẫu số kia:
\begin{itemize}\setlength{\itemsep}{1pt}\setlength{\leftskip}{-2.5mm}
  \item Chọn mẫu số lớn hơn làm \textbf{MSC}.
  \item Lấy MSC chia cho mẫu số bé để tìm \textbf{thừa số phụ}: $k = \text{MSC} : b$.
  \item Nhân cả tử và mẫu của phân số có mẫu bé với $k$:
  $$\dfrac{a}{b} = \dfrac{a \times k}{b \times k}$$
  \item \textbf{Giữ nguyên} phân số có mẫu số lớn!
\end{itemize}
}
\end{tcolorbox}
\end{minipage}

\vspace{-2.5mm}
% --- KHỐI 3: SƠ ĐỒ TƯ DUY TIKZ ---
\noindent
\begin{tcolorbox}[
  enhanced,
  colback=orange!4!white,
  colframe=orange!80!black,
  arc=2mm, boxrule=0.8pt,
  left=2mm, right=2mm, top=1.5mm, bottom=1.5mm,
  title={\footnotesize\bfseries\color{white}\faSitemap\ SƠ ĐỒ TƯ DUY: 2 NHÁNH QUY ĐỒNG MẪU SỐ TỐI ƯU}
]
\centering
\begin{tikzpicture}[
  box/.style={draw, rounded corners=2pt, align=center, font=\scriptsize\bfseries, inner sep=3pt, minimum height=1.1cm},
  arr/.style={-{Stealth[length=2.5mm]}, thick}
]
  \node[box, fill=purple!12, draw=purple!80!black, text width=3.2cm] (root) at (0, 0) {QUY ĐỒNG\\MẪU SỐ\\$\dfrac{a}{b}$ và $\dfrac{c}{d}$};

  \node[box, fill=blue!10, draw=blue!80!black, text width=4.8cm] (b1) at (5.8, 0.95) {CÁCH TỔNG QUÁT (Nhân chéo)\\$\text{MSC} = b \times d$\\Nhân tử \& mẫu phân số này với mẫu kia};

  \node[box, fill=teal!10, draw=teal!80!black, text width=4.8cm] (b2) at (5.8, -0.95) {MẪU LỚN CHIA HẾT MẪU BÉ\\$\text{MSC} =$ Mẫu lớn\\Nhân mẫu bé với thừa số phụ, GIỮ NGUYÊN mẫu lớn!};

  \draw[arr, blue!80!black] (root.east) -- ++(0.5, 0) |- (b1.west);
  \draw[arr, teal!80!black] (root.east) -- ++(0.5, 0) |- (b2.west);
\end{tikzpicture}
\end{tcolorbox}

\vspace{-2.5mm}
% --- KHỐI 4: VÍ DỤ & MẸO TRÁNH BẪY ---
\noindent
\begin{minipage}[t]{0.49\textwidth}\vspace{0pt}
\begin{tcolorbox}[
  enhanced,
  colback=green!4!white,
  colframe=green!60!black,
  arc=2mm, boxrule=0.8pt,
  left=2mm, right=2mm, top=1.8mm, bottom=1.8mm,
  title={\footnotesize\bfseries\color{white}\faCheckCircle\ Ví dụ thực hành chuẩn mực}
]
{\footnotesize
\textbf{1. Quy đồng $\dfrac{2}{3}$ và $\dfrac{4}{5}$ (MSC = 15):}\\
$\dfrac{2}{3} = \dfrac{2 \times 5}{3 \times 5} = \mathbf{\dfrac{10}{15}}; \quad \dfrac{4}{5} = \dfrac{4 \times 3}{5 \times 3} = \mathbf{\dfrac{12}{15}}$.\\
\textbf{2. Quy đồng $\dfrac{3}{4}$ và $\dfrac{5}{12}$ (MSC = 12):}\\
Vì $12 : 4 = 3$ (thừa số phụ là 3):\\
$\dfrac{3}{4} = \dfrac{3 \times 3}{4 \times 3} = \mathbf{\dfrac{9}{12}}$; \quad Giữ nguyên $\mathbf{\dfrac{5}{12}}$.
}
\end{tcolorbox}
\end{minipage}
\hfill
\begin{minipage}[t]{0.49\textwidth}\vspace{0pt}
\begin{tcolorbox}[
  enhanced,
  colback=red!4!white,
  colframe=red!70!black,
  arc=2mm, boxrule=0.8pt,
  left=2mm, right=2mm, top=1.8mm, bottom=1.8mm,
  title={\footnotesize\bfseries\color{white}\faExclamationTriangle\ Mẹo tránh bẫy trắc nghiệm}
]
{\footnotesize
\begin{itemize}\setlength{\itemsep}{1.5pt}\setlength{\leftskip}{-2.5mm}
  \item \textbf{Bẫy lạm dụng nhân chéo:}\\
  Với $\frac{3}{4}$ và $\frac{5}{12}$, nếu nhân chéo $\text{MSC} = 48$ thì số rất to, tính toán dễ nhầm. Luôn kiểm tra xem mẫu lớn có chia hết cho mẫu bé không!
  \item \textbf{Bẫy quên nhân tử số:}\\
  Khi nhân mẫu với thừa số phụ, học sinh hay quên nhân cả tử số.
\end{itemize}
}
\end{tcolorbox}
\end{minipage}

\end{document}
